\chapter[RESULTADOS Y DISCUSIÓN]{RESULTADOS Y DISCUSIÓN}
\label{cap:resultados}


\begin{borrador}
  \item Capítulo vivo mientras se hacen las fases, a escribir mientras pasan las tarjetas y fases.
  \item \textbf{Regla de este capítulo, que conviene tener delante al escribir:} aquí van los
  resultados, y en la sección de discusión su interpretación. Si una frase de las secciones
   anteriores empieza por <<esto implica>> o <<esto demuestra>>, pertenece a la discusión. Es
  el error más común en este capítulo y el más fácil de cometer.
  \item \textbf{Orden de lectura dentro de cada sección}, que también es convención y ayuda
  al lector: qué se midió y sobre qué colección; la tabla o figura; el resultado principal;
  los desgloses; y las lecturas erróneas que conviene desactivar antes de que las haga otro.
  \item \sout{falta texto introductorio del capítulo.} \textbf{Debajo lo tienes en
    turquesa}; reescríbelo y quítale el color. Si cambia el reparto de secciones, el primer
  párrafo deja de ser cierto, porque anuncia el orden.
\end{borrador}

\textcolor{borrador2}{Este capítulo recoge lo que el sistema produce y lo que se ha medido
  sobre él. Se organiza siguiendo los objetivos específicos enunciados en el
  Capítulo~\ref{cap:objetivos}: la colección documental que se construye, la fragmentación
  con la que se parte, el modelo de incrustaciones y la base de datos vectorial con los que
  se representa y se indexa, la evaluación del sistema completo y su validación con usuarios.
  La última sección discute qué responde cada resultado a cada objetivo y con qué alcance.}

\textcolor{borrador2}{Todas las cifras proceden de la ejecución de los verificadores y de los
  guiones de experimentación sobre la colección completa del curso 2026-27, y no sobre una
  muestra. El método con el que se obtuvieron ---qué se compara, con qué instrumento y por
  qué--- está en el Capítulo~\ref{cap:materialesymetodos}; aquí están las cifras y su
  lectura.}
\section[La colección documental construida]{La colección documental construida} 
\label{secc:resultados_coleccion}

La colección documental o \textit{corpus} resultante recoge las 12 titulaciones disponibles en la EPSJ en el curso 2026-27, siete grados y cinco dobles grados, con sus 528 asignaturas, las 288 guías docentes publicadas y 8 bloques de salidas profesionales.

De esas 528 asignaturas, 442 disponen de guía docente, una cobertura del 83,7\%. Las 86 restantes no la tienen ya que no se han publicado aún a fecha de este trabajo.\sout{,} El reparto de estas guías faltantes no es uniforme, se concentran en las titulaciones nuevas que todavía no se han completado. Conviene tener este dato en cuenta al leer cualquier resultado posterior en esta memoria, porque significa que el sistema puede saber menos precisamente de las titulaciones más nuevas.

Sobre esta recopilación de información, el fragmentador produce 1.334 fragmentos repartidos en \sout{332} \textcolor{borrador2}{322} unidades, de las cuales 78 se comparten entre varias titulaciones al repetirse en los diferentes grados. Los tipos son: 1.193 fragmentos de guías docentes, 86 nombres de asignaturas sin guía disponible, 33 de plan de estudios y 22 salidas profesionales. Los tamaños de estos fragmentos van desde los 171 a 900 caracteres, con una mediana de 838 y un percentil 90 de 894. Ninguno supera la ventana del modelo de incrustaciones.



\textcolor{borrador2}{Estas cifras no proceden de una fuente uniforme, y conviene
  dejar constancia de ello porque condiciona lo que la colección puede llegar a contener. La
  web de la EPSJ no publica ningún esquema de sus páginas: cada titulación las genera con la
  plantilla de su plan y esas plantillas no coinciden entre sí. El repositorio recoge cinco
  registros de anomalías de calidad de datos ---la serie DQA--- con la evidencia de cada
  una\cite{repo_tfg}, y todas comparten un rasgo: casi ninguna produce un error.}

\textcolor{borrador2}{La primera afecta al formato. Las 288 guías docentes del curso 2026-27
  se sirven como PDF y no como HTML, detrás de una dirección que sigue acabando en
  \texttt{.html}, de modo que el formato solo se conoce mirando el tipo de contenido de la
  respuesta. El PDF incluye además un bloque de profesorado con nombre, categoría, despacho,
  correo electrónico y teléfono ---entre cuatro y seis correos y entre dos y cuatro teléfonos
  en los ejemplares analizados--- que la versión en HTML no traía. Como la colección excluye
  el profesorado a propósito, la extracción admite únicamente las secciones de resumen y de
  descripción de contenidos y descarta por defecto todo lo demás.}

\textcolor{borrador2}{La segunda afecta a la estructura de las tablas de asignaturas, que no
  es la misma en todos los planes. El plan de 2025 del Grado en Ingeniería Geomática y
  Topográfica intercala una columna de curso recomendado en las tablas de mención, que reúnen
  19 de las 39 asignaturas de la titulación; y los dobles grados rotulan «carácter» la columna
  que los grados simples rotulan «tipo», y abrevian «OBL» donde los simples escriben «OB».
  Leída por la posición de la columna, una tabla así no falla: devuelve cada dato de la
  columna equivocada y la fila se descarta sin excepción y sin traza. Por eso las columnas se
  localizan por su rótulo.}

\textcolor{borrador2}{La tercera afecta a los dobles grados, que son cinco de las doce
  titulaciones. Publican su plan de estudios bajo una ruta distinta, escriben los nombres en
  mayúsculas y a veces con el acrónimo de la titulación de procedencia, y usan una serie de
  códigos propia: ninguno de los códigos de sus asignaturas coincide con el de la misma
  asignatura en el grado base, de manera que la correspondencia entre unas y otras solo puede
  establecerse por el nombre.}

\textcolor{borrador2}{La cuarta afecta a lo que las guías dicen. De las 288 publicadas, cinco
  traen la plantilla completa ---sus trece rótulos en su sitio--- y las secciones de resumen y
  de contenidos vacías desde el origen. No son guías ilegibles ni mal extraídas: están
  publicadas y no dicen nada, y todas pertenecen a planes en implantación. La colección las
  distingue de las asignaturas cuya guía no se ha publicado todavía, porque no son lo mismo:
  una existe y está vacía, la otra no existe.}

\textcolor{borrador2}{Quedan las irregularidades menores, que no cambian las cifras pero sí
  el código que hace falta para obtenerlas: direcciones que no siguen ningún patrón deducible
  ---construirlas produce errores 404 reales---, tablas que mezclan asignaturas troncales con
  optativas agrupadas por mención, el Trabajo Fin de Grado declarado con tipo propio en unos
  planes y como obligatoria en otros, filas de relleno sin código que no corresponden a
  ninguna asignatura, espacios duros y caracteres de ancho cero intercalados en el texto,
  direcciones con la extensión repetida, asteriscos de nota al pie pegados al nombre y la
  marca de asignatura no ofertada dentro del propio nombre. Cada una se trata con una regla
  determinista, y ninguna se resuelve inventando el dato: el crédito que la fuente no publica
  se queda ausente en la colección.}

\begin{borrador}
  \item \textbf{Andamiaje en turquesa: la prosa la escribes tú.} Son las anomalías de la
  fuente que condicionan lo que la colección puede contener, resumidas de los cinco registros
  DQA del repositorio.
  \item \textbf{Cada afirmación tiene su registro detrás}, por si te la discuten: el formato
  en PDF y el bloque de profesorado, el DQA-0002; las tablas con las columnas cambiadas, el
  DQA-0003; los dobles grados, el DQA-0005; las guías publicadas vacías, el DQA-0004; y las
  irregularidades menores, el DQA-0001.
  \item \textbf{Decide si los citas uno a uno o solo el repositorio.} Citarlos por su
  identificador ayuda a quien quiera comprobarlo, pero obliga a que los cinco estén
  accesibles en la entrega.
  \item \textbf{El hilo que las une, que es lo que hay que poder decir en una frase:} casi
  ninguna de estas irregularidades produce un error. Una tabla con las columnas cambiadas se
  lee sin excepción y devuelve datos de la columna equivocada, y una titulación entera
  desaparece en silencio. De ahí que la colección se compruebe con verificadores propios sobre
  el conjunto completo y no solo con las pruebas del código.
\end{borrador}

\section[Fragmentación de la colección de datos]{Fragmentación de la colección de datos}
\label{secc:resultados_fragmentacion}


El experimento de rejilla descrito en el Apartado~\ref{secc:rejilla_fragmentacion} produce cuarenta y cinco configuraciones, todas ellas explicadas en el Registro de Decisiones de Arquitectura \mbox{ADR-0001}\cite{repo_tfg}. 

\sout{Cada uno de los resultados se presentan separando dos factores} \textcolor{borrador2}{Cada uno de los resultados se presenta separando dos factores}, que es por lo que la rejilla se diseñó de manera simétrica. Primero se pregunta si es que la estrategia en sí de corte importa, y para responder eso, hay que elegir un tamaño y dejarlo fijo y comparar solo los diferentes cortes. Después de esto, se pregunta cuánto importa el tamaño, y para ello hay que dejar fija la estrategia y de nuevo ver los diferentes tamaños.

\subsection[Efecto de la estrategia de corte]{Efecto de la estrategia de corte}
\label{secc:rejilla_estrategia}

La Tabla~\ref{tab:rejilla_fragmentacion} recoge la mejor configuración
de las tres estrategias. Las cuarenta y cinco están en el ADR-0001\cite{repo_tfg}.

\begin{table}[htbp]
  \centering
  \caption[Efecto de la estrategia y del tamaño en la fragmentación]{\textcolor{borrador2}{Rejilla
      de fragmentación, en dos bloques. Arriba, las tres estrategias a igualdad de tamaño
      máximo; abajo, el efecto del tamaño con la estrategia fija. Cincuenta preguntas del
      conjunto de evaluación.}}
  \label{tab:rejilla_fragmentacion}
  \small
  \begin{tabular}{@{}llrrrrrr@{}}
    \toprule
    \textbf{Estrategia}                                                                              & \textbf{Ajuste}         & \textbf{Máx.}  & \textbf{Frag.} & \textbf{RU@1}  & \textbf{RU@3}  & \textbf{MRR}   & \textbf{Trunc.} \\
    \midrule
    \multicolumn{8}{@{}l}{\textit{(a) Las tres estrategias, con el tamaño máximo fijado en 900}}                                                                                                                              \\
    \addlinespace[2pt]
    Tamaño fijo                                                                                      & solape 20 \%            & 900            & 1.477          & \textbf{0,935} & 0,965          & \textbf{0,974} & 0               \\
    \textbf{Estructural}                                                                             & objetivo 100 \%         & 900            & 1.334          & 0,930          & \textbf{0,985} & 0,970          & 0               \\
    Tamaño fijo                                                                                      & solape 0 \%             & 900            & \textbf{1.228} & 0,930          & 0,985          & 0,963          & 0               \\
    Semántica                                                                                        & percentil 70            & 900            & 1.647          & 0,930          & 0,985          & 0,963          & 0               \\
    \midrule
    \multicolumn{8}{@{}l}{\textit{(b) El tamaño máximo, con la estrategia fijada en tamaño fijo sin solape}}                                                                                                                  \\
    \addlinespace[2pt]
    Tamaño fijo                                                                                      & solape 0 \%             & 600            & 2.312          & 0,950          & 0,985          & 0,972          & 0               \\
    Tamaño fijo                                                                                      & solape 0 \%             & 900            & 1.228          & 0,930          & 0,985          & 0,963          & 0               \\
    Tamaño fijo                                                                                      & solape 0 \%             & 1.200          & 912            & 0,910          & 0,985          & 0,960          & 0               \\
    Tamaño fijo                                                                                      & solape 0 \%             & 1.500          & 736            & 0,850          & 0,985          & 0,927          & 3               \\
    Tamaño fijo                                                                                      & solape 0 \%             & 1.800          & 623            & 0,790          & 0,945          & 0,885          & 20              \\
    \midrule
    \multicolumn{8}{@{}l}{\textit{(c) La configuración que empleaba el sistema antes de esta comparación}}                                                                                                                    \\
    \addlinespace[2pt]
    \textit{Estructural}                                                                             & \textit{objetivo 80 \%} & \textit{1.500} & \textit{884}   & \textit{0,780} & \textit{0,965} & \textit{0,875} & \textit{0}      \\
    \bottomrule
  \end{tabular}
  \begin{minipage}{0.95\textwidth}
    \vspace{4pt}
    \tiny
    \textbf{Notas:} \textbf{Máx.}: tamaño máximo del fragmento, en caracteres;
    \textbf{Frag.}: número total de fragmentos que produce esa configuración sobre la
    colección completa; \textbf{RU@K} (\textit{Recall por Unidad}): proporción de preguntas
    cuya asignatura correcta aparece entre los $K$ primeros resultados; \textbf{MRR}
    (\textit{Mean Reciprocal Rank}): media del inverso de la posición en la que aparece el
    primer resultado correcto, de modo que premia colocarlo antes; \textbf{Trunc.}: número
    de fragmentos que superan la ventana del modelo de incrustaciones y que este recorta en
    silencio, sin avisar ni fallar. En negrita, el mejor valor de cada columna dentro del
    bloque (a). \textcolor{borrador2}{Una diferencia de 0,020 en RU@1 equivale a una
      pregunta de las cincuenta.}
  \end{minipage}
\end{table}


La tabla de arriba contesta a \sout{la duda que estamos abarcando esta subsección} \textcolor{borrador2}{la duda que aborda esta subsección}, y es que a igualdad de tamaño, las tres estrategias son prácticamente indistinguibles, dando resultados muy similares. Las nueve configuraciones diferentes (3 estructural, 3 de corte por longitud y 3 de corte semántico), que comparten el máximo de 900 caracteres, se reparten entre 0,890 y 0,935 de acierto en el primer resultado, valores que no son de vital importancia.

Hay que aclarar que la recuperación no justifica por sí sola la estrategia elegida, ya que en el caso del troceado por longitud sin solape, con este mismo máximo, produce ciento seis fragmentos menos que la estructural y empata con ella en casi todas las medidas de acierto, \sout{por lo que si solo nos moviéramos por el rendimiento medido, eligiríamos esta sin dudarlo} \textcolor{borrador2}{de modo que, si solo contara el rendimiento medido, se elegiría esta sin dudarlo}.

Lo que separa una estrategia de otra no aparece en ningún valor medido y es el de la metodología, y es que la estrategia estructural parte el texto por sus propias fronteras, mientras que la de longitud fija parte al alcanzar un número máximo de caracteres, con lo que empieza y termina a mitad de una frase y a veces a mitad de una palabra. Estas características para las medidas del experimento dan igual y dan incluso mejores resultados, pero esta información se le va a pasar a un modelo generativo del lenguaje para que redacte la respuesta y ahí una frase truncada a la mitad sí que importa, \sout{puesto que no le estaríamos dando al usuario toda la información} \textcolor{borrador2}{puesto que no se le estaría dando al usuario toda la información} disponible de ese apartado.




\subsection[Efecto del tamaño del fragmento]{Efecto del tamaño del fragmento}
\label{secc:rejilla_tamano}

La segunda pregunta (cuánto importa el tamaño) se responde con la estrategia ya fijada, en este caso la estructural, de modo que lo único que varía entre las filas de la Tabla~\ref{tab:tamano_fragmento} sea el máximo. 


% !TEX root = ../main.tex
\begin{table}[htbp]
  \centering
  \caption[Efecto del tamaño máximo con la estrategia elegida]{Efecto
      del tamaño máximo del fragmento, manteniendo fija la estrategia estructural con el
      objetivo igualado al máximo.}
  \label{tab:tamano_fragmento}
  \small
  \begin{tabular}{@{}rrrrrrr@{}}
    \toprule
    \textbf{Máx.} & \textbf{Frag.} & \textbf{Mediana} & \textbf{RU@1}  & \textbf{RU@3}  & \textbf{MRR}   & \textbf{Trunc.} \\
    \midrule
    600           & 2.695          & 569              & 0,930          & 0,980          & 0,960          & 0               \\
    \textbf{900}  & \textbf{1.334} & \textbf{838}     & \textbf{0,930} & \textbf{0,985} & \textbf{0,970} & \textbf{0}      \\
    1.200         & 947            & 1.104            & 0,850          & 0,960          & 0,913          & 0               \\
    1.500         & 753            & 1.344            & 0,830          & 0,945          & 0,898          & 4               \\
    1.800         & 632            & 1.497            & 0,770          & 0,945          & 0,868          & 11              \\
    \bottomrule
  \end{tabular}
  \begin{minipage}{0.95\textwidth}
    \vspace{4pt}
    \tiny
    \textbf{Notas:} \textbf{Máx.}: tamaño máximo del fragmento, en caracteres, que es una
    restricción dura; \textbf{Frag.}: número total de fragmentos que produce esa
    configuración sobre la colección completa; \textbf{Mediana}: tamaño mediano del
    fragmento resultante, en caracteres; \textbf{RU@K} (\textit{Recall por Unidad}):
    proporción de preguntas cuya asignatura correcta aparece entre los $K$ primeros
    resultados; \textbf{MRR} (\textit{Mean Reciprocal Rank}): media del inverso de la
    posición del primer resultado correcto; \textbf{Trunc.}: fragmentos que superan la
    ventana del modelo de incrustaciones, contados con su propio analizador léxico, y que el
    modelo recorta en silencio, sin avisar ni fallar. En negrita, la configuración adoptada.
  \end{minipage}
\end{table}


El valor 900 es el único que gana en las tres métricas a la vez, y lo consigue con la mitad de fragmentos que el de 600 caracteres. En el documento que recoge el experimento (\mbox{ADR-0001}) hay una cosa que se repite constantemente, y es que a más fragmentos, mejores métricas, ya que en un corte más fino, la probabilidad de pedir un fragmento y que el recogido sea ese exactamente son más altas.

La última columna de esta tabla aporta un argumento que dejaba instantáneamente las configuraciones \sout{con máximos de 1500 y 1800} \textcolor{borrador2}{con máximos de 1.500 y 1.800}, y es que estos truncaban los fragmentos y al hacer esto, mucha información se perdía en el camino, algo que es inadmisible.

En resumen de este apartado, la configuración es el corte estructural con un tamaño máximo de 900 y mínimo de 200 caracteres. Por encima de 1.200 caracteres el modelo empieza a recortar fragmentos, y de las opciones por debajo de esta cifra, el de 900 es el único en ganar en tres categorías, con la mitad de fragmentos que el de 600. Entre tres opciones con métricas prácticamente iguales, se conserva la que corta por las fronteras de texto, cuya ventaja está más que en la recuperación, en la generación.


\begin{borrador}
  \item \textbf{Objeción de defensa:} «si el modelo de incrustaciones ya estaba elegido y se
  conocía su ventana, ¿por qué se probaron tamaños máximos que iban a truncar?».
  \item \textbf{Primero, por qué hacía falta un modelo.} El experimento mide recuperación, y
  recuperar exige vectores: no hay manera de comparar troceados por su acierto sin un modelo
  de incrustaciones detrás. Eso crea un bucle ---para fijar el troceado hace falta el modelo, y
  la comparación de modelos necesitó a su vez un troceado--- que solo se rompe fijando uno de
  los dos. Aquí se fija el modelo, y se fija el que el sistema usa de verdad, con lo que el
  troceado queda ajustado al modelo real y no a uno de referencia. El precio de romperlo así
  está declarado entre las amenazas: el espacio conjunto de las dos decisiones no se ha
  recorrido.
  \item \textbf{Y segundo, por qué se probaron tamaños que truncan.} Porque que 900 caracteres
  quepan en la ventana es un \emph{resultado} del experimento y no un supuesto de partida. La
  equivalencia entre caracteres y \textit{tokens} no es fija: los números, las abreviaturas,
  los códigos de asignatura y las palabras poco frecuentes consumen más \textit{tokens} por
  carácter que la prosa corriente, así que un máximo en caracteres solo se traduce a un máximo
  en \textit{tokens} midiéndolo sobre el corpus real. Los tamaños grandes son el control que
  localiza dónde empieza el recorte en lugar de suponerlo.
\end{borrador}




\begin{borrador}
  \item \textbf{Falta un aviso que evita una contradicción aparente:} el troceo por
  longitud de esta tabla \emph{no} es el que descarta la
  Tabla~\ref{tab:alternativas_chunking}. Allí se descarta uno que recorre la colección entera
  y puede juntar dos asignaturas en un fragmento; aquí las tres estrategias trabajan ya dentro
  de una unidad. Sin esa frase, una tabla lo descarta y la otra lo pone arriba.
\end{borrador}


\begin{borrador}                                                                                                                                                    
  \item \textbf{Fase:} 1 (IT-16). Aquí van las cifras y su lectura; el diseño de la rejilla
  ---qué se comparó y con qué instrumento--- se queda en el
  Capítulo~\ref{cap:materialesymetodos}. Si al releer encuentras aquí una frase que explica
  \emph{cómo} se midió, es que pertenece allí.
  \item \textbf{Cuidado con cómo lo cuentas}, que aquí es fácil pasarse de listo: la
  conclusión no es «la fragmentación semántica es peor». Es que \textbf{no es mejor}, que es
  distinto y más difícil de rebatir. Se descarta por coste, no por calidad.
  \item \textbf{Amenaza que hay que enlazar} con la Sección~\ref{secc:amenazas_validez}: las
  diferencias de una pregunta sobre cincuenta ---0,020--- no son distinguibles de lo que
  movería otra anotación. La decisión se apoya en la comparación a igualdad de fragmentos y
  en el recuento de truncados, no en el orden bruto de la tabla.
  \item \textbf{Respuesta preparada, por si sale:} «mide usted los dos efectos por
  separado, pero podrían depender el uno del otro». Dependen, y está medido: la distancia
  entre estrategias vale 0,000 con 600 caracteres y 0,050 con 1.800. Lo que se sostiene es que
  no dependen por debajo de 1.200, que es donde se decide, y que por encima la decisión ya la
  ha tomado el truncado. Las quince medias que lo sostienen salen del anexo del ADR-0001.
  \item \textbf{Lo que no se ha medido:} el efecto sobre la generación, que es donde la
  fragmentación actúa de verdad. Si un fragmento parte una definición y solo llega la mitad,
  la recuperación puntúa igual de bien y la respuesta sale peor. Exige métricas de la Fase 2.
\end{borrador}

\section{Selección del modelo de incrustaciones}
\label{secc:resultados_embeddings}

\begin{borrador}
  \item \textbf{LA TABLA DE ESTA SECCIÓN ESTÁ CADUCADA (IT-16).} Sus cifras son de la
  ejecución del 04/08/2026 sobre 797 fragmentos, con el máximo de fragmento todavía en
  1.500 caracteres. Al fijar los parámetros de fragmentación el corpus pasó a
  \textbf{1.334} y la comparativa se repitió el 06/08/2026. Las cifras vigentes están en
  \texttt{docs/experimentos/it28-embeddings.md}, que escribe el propio guion: cópialas de
  ahí a \texttt{tablas/TablaComparativaEmbeddings.tex}, no a mano desde aquí.
  \item \textbf{Y hay que reescribir la conclusión, no solo los números.} El texto de
  abajo dice que el modelo grande recupera mejor y que se descarta por coste. Con el
  troceado nuevo la diferencia entre los dos baja de dos preguntas y media de cincuenta a
  \textbf{media pregunta}, y el MRR se iguala en 0,970. La elección no cambia; el
  argumento sí, y a mejor: ya no depende de que el sistema acabe usando diez resultados
  por consulta, que es un parámetro que todavía no está decidido.
  \item \textbf{Se da la vuelta además el desglose por tipo de pregunta.} En las de
  metadatos ---donde el modelo pequeño era claramente peor--- pasa a ir por delante
  (0,697 frente a 0,664). Si dejas la frase anterior sin tocar, estarás defendiendo la
  decisión con una razón que hoy es falsa.
  \item \textbf{Lo que no se puede decir:} que el modelo pequeño haya mejorado. El modelo
  es el mismo y lo que cambió fue la colección. Trocear más fino reparte cada unidad en
  más fragmentos y le da más oportunidades de entrar entre los primeros resultados, y de
  eso se beneficia más el modelo que peor las aprovechaba.
  \item \textbf{Cuidado con la exhaustividad por fragmento, que baja} (0,771 a 0,697) y
  \emph{no} es una regresión: su techo baja con ella, de 0,905 a 0,789. Medido contra el
  techo, el hueco se reduce de 0,134 a 0,092. Sin el techo al lado, la cifra dice lo
  contrario de lo que ocurre.
\end{borrador}

\begin{borrador}      
  \item \textbf{Fase: 1 (IT-28, IT-29, IT-100).} Sección nueva, con los resultados del
  experimento del 04/08/2026. Andamiaje: la prosa la escribes tú. La tabla ya está
  puesta y sus cifras salen del fichero que genera el propio guion, así que
  \textbf{no las copies a mano en el texto}: si cambian, cambian ahí.
  \item Regla del capítulo, para no repetir el error de otras secciones: aquí van
  \emph{los resultados}, no la justificación del método. El diseño del experimento y
  el porqué de cada candidato es Capítulo~\ref{cap:materialesymetodos}. Y la
  interpretación fina («qué significa esto para el sistema») no debería mezclarse con
  el dato.
  \item Orden sugerido para redactar: (1) qué se midió y sobre qué; (2) la tabla;
  (3) el resultado del modelo en español, que es el hallazgo; (4) el desglose por tipo
  de pregunta; (5) las dos lecturas erróneas que hay que evitar.
  \item Lo que \textbf{no} se ha medido y hay que decir en algún sitio: la latencia de
  una consulta, ni con este modelo ni con ningún otro. Los tiempos que se manejan son
  de incrustar la colección entera, que es una operación por lotes y poco frecuente.
\end{borrador}

\textcolor{borrador2}{La comparativa se ejecutó sobre la colección completa, con las
  cincuenta preguntas anotadas del conjunto de evaluación, y sus resultados se recogen en la
  Tabla~\ref{tab:comparativa_embeddings}. Los cuatro candidatos admiten al menos quinientos
  doce \glr{tokens}{token}, de modo que ninguno se queda sin leer parte de los fragmentos:
  esa condición se comprueba en cada ejecución y no se da por supuesta, porque en las
  ejecuciones anteriores del experimento dos de los modelos evaluados solo alcanzaban a leer
  la mitad del texto y sus cifras no eran comparables con las del resto.}

% !TEX root = ../main.tex
% Resultados literales de la ejecucion del 06/08/2026.
% Fuente: anexo del ADR-0003, escrito por el propio script.
% NO editar estas cifras a mano: si cambian, se vuelve a ejecutar
%   py scripts/experimentos/experimento_embeddings.py
% y se copian de ahi. Corpus de 1.334 fragmentos y 50 preguntas.
\begin{table}[htbp]
  \centering
  \small
  \caption[Comparativa de modelos de incrustaciones]{Comparativa de modelos de
    incrustaciones sobre la colección completa y las 50 preguntas del
    conjunto de evaluación.
    Fuente: elaboración propia.}
  \label{tab:comparativa_embeddings}
  % Ocho columnas y un nombre de modelo largo en monoespaciada desbordaban la
  % caja 14,6 pt. Se aprieta la separacion entre columnas en lugar de encoger
  % la letra: 2 pt menos a cada lado de ocho columnas son 32 pt, de sobra.
  \setlength{\tabcolsep}{4pt}
  \begin{tabular}{@{}lrrrrrrr@{}}
    \toprule
    \textbf{Modelo}                            & \textbf{R@3}   & \textbf{R@5}   & \textbf{R@10}
                                               & \textbf{RU@3}  & \textbf{RU@5}  & \textbf{RU@10} & \textbf{MRR}   \\
    \midrule
    \texttt{multilingual-e5-small}             & 0,697          & 0,803          & 0,911
                                               & 0,985          & 0,990          & \textbf{1,000} & \textbf{0,970} \\
    \texttt{multilingual-e5-large}             & \textbf{0,740} & \textbf{0,852} & \textbf{0,938}
                                               & \textbf{0,995} & \textbf{1,000} & \textbf{1,000} & \textbf{0,970} \\
    \texttt{bge-m3}                            & 0,720          & 0,836          & 0,917
                                               & 0,985          & 0,995          & 0,995          & 0,949          \\
    \texttt{sentence\_similarity\_spanish\_es} & 0,152          & 0,194          & 0,307
                                               & 0,410          & 0,505          & 0,610          & 0,342          \\
    \bottomrule
  \end{tabular}
  \begin{minipage}{0.95\textwidth}
    \vspace{4pt}
    \tiny
    \textbf{Notas:} \textbf{R@K} es la exhaustividad por \emph{fragmento} (cuántos de los
    trozos de la unidad correcta aparecen entre los $K$ primeros resultados).
    \textbf{RU@K} es la exhaustividad por \emph{unidad} (si se ha encontrado la asignatura
    correcta, sin penalizar que falte alguno de sus trozos). \textbf{MRR}
    (\textit{Mean Reciprocal Rank}) es la media del inverso de la posición del primer resultado
    correcto, de modo que premia colocarlo antes. En negrita, el mejor valor de cada
    columna. El
    modelo adoptado es \texttt{multilingual-e5-small}.
  \end{minipage}
\end{table}


\textcolor{borrador2}{El resultado que más información aporta no es cuál de los modelos
  obtiene la mejor puntuación, sino la distancia a la que queda el único candidato entrenado
  específicamente en español. Su exhaustividad por unidad en los tres primeros resultados es
  de 0,320 frente al 0,945 del modelo multilingüe de tamaño equivalente, y no recupera ninguna
  de las catorce preguntas que piden el listado completo de asignaturas de una titulación. Al
  haber leído exactamente el mismo texto que los demás, la diferencia no puede atribuirse a
  que dispusiera de menos información, sino a la tarea para la que fue entrenado: reconocer si
  dos frases significan lo mismo, que no es lo que se le pide a un sistema de recuperación,
  donde una pregunta breve debe emparejarse con un documento extenso que la responde y que no
  se le parece. Conviene precisar el alcance de esta conclusión: no permite afirmar que no
  pueda existir un modelo en español adecuado para recuperación, sino que el mejor disponible
  en el momento de la comparación no lo es.}

\begin{borrador}      
  \item \textbf{Dos lecturas erróneas de la tabla que conviene desactivar tú mismo}, porque
  son las dos preguntas que va a hacer el tribunal.
  \item \textbf{«Un 0,697 significa que falla una de cada tres veces.»} No. La
  exhaustividad por fragmento \textbf{no tiene techo 1}: una unidad repartida en más de $K$
  fragmentos no cabe entera en los $K$ primeros resultados. Sobre esta colección el máximo
  alcanzable es \textbf{0,789} para K=3, \textbf{0,906} para K=5 y \textbf{0,968} para K=10.
  Hay que restar del techo, no de la unidad.
  \item \textbf{«\texttt{bge-m3} recupera mejor que el modelo elegido.»} Solo en una columna,
  y no es la que decide. Gana en exhaustividad \emph{por fragmento} (0,720 frente a 0,697 con
  K=3), empata en la \emph{por unidad} con K=3, y con K=10 la pierde: 0,995 frente al
  \textbf{1,000} del pequeño. Y en MRR queda claramente por detrás ---0,949 frente a
  0,970---, es decir, coloca el primer acierto más abajo.
  \item \textbf{Y esa ventaja no está repartida, está concentrada.} En el desglose por tipo,
  \texttt{bge-m3} gana solo en las preguntas de salidas profesionales (0,917 frente a 0,690);
  en las de metadatos \textbf{es el peor de los tres} (0,631 frente a 0,697 del pequeño). Una
  media general que no se desglosa aquí diría lo contrario de lo que ocurre.
  \item \textbf{Cuidado con explicar por qué.} \texttt{bge-m3} se diferencia del modelo
  elegido en tamaño, familia, entrenamiento y convención de prefijos \emph{a la vez}, así que
  este experimento \textbf{no puede atribuir} la diferencia a ninguna de esas causas. El único
  par que aísla una sola variable es el de los dos modelos de la misma familia. Decirlo así es
  más fuerte que aventurar una explicación.
  \item \textbf{Y una tercera, sobre el crecimiento con K.} La exhaustividad es monótona
  creciente en K, así que que K=10 supere a K=5 es una propiedad de la métrica y no un
  hallazgo. Lo que decide K es el coste de contexto del generador, y eso se mide en la
  Fase 2.
\end{borrador}

\section[Base de datos vectorial]{\textcolor{borrador2}{Base de datos vectorial}}
\label{secc:resultados_vectordb}

\begin{borrador}
  \item \textbf{Sección nueva y vacía (07/08/2026).} Completa la otra mitad del \textbf{OE3},
  que pide elegir con criterios cuantitativos el modelo de incrustaciones \emph{y} la base de
  datos vectorial. Hoy solo está la primera mitad.
  \item Pendiente de IT-31 e IT-32. A diferencia de la comparación de modelos, esta se
  diseñó \textbf{con los umbrales escritos antes de ejecutar nada}, que es justamente lo que
  corrige la amenaza principal del ADR-0003: conviene decirlo con esas palabras, porque
  demuestra que la lección se aplicó.
  \item Y hay que declarar aquí una amenaza que no cabe en otro sitio: con menos de mil
  quinientos fragmentos, \textbf{la búsqueda exacta por fuerza bruta es competitiva}, de modo
  que la base vectorial no se puede justificar por velocidad sin haberla medido.
\end{borrador}

\section[Evaluación del \textit{pipeline} RAG]{\textcolor{borrador2}{Evaluación del \textit{pipeline} RAG}}
\label{secc:resultados_rag}

\begin{borrador}
  \item \textbf{Renombrada el 07/08/2026:} se llamaba <<Evaluación de sistemas RAG>>, que
  suena a repaso del estado del arte y no a resultado propio. Aquí van las métricas de
  \emph{este} sistema, que es lo que responde al \textbf{OE5}.
\end{borrador}

\section{Validación con usuarios}

\begin{borrador}      
  \item \textbf{Fase:} 4 Validación de usuarios reales (los tengo que buscar aún y no sé si será posible)
\end{borrador}

\section[Discusión]{\textcolor{borrador2}{Discusión}}
\label{secc:discusion}

\begin{borrador}
  \item \textbf{Sección nueva, y esta no es opcional (07/08/2026).} La normativa de la EPSJ
  nombra el capítulo <<Resultados \textbf{y Discusión}>> como una sola unidad, y entre los
  criterios con los que se evalúa el trabajo está la \emph{correspondencia entre los
    objetivos de la propuesta original y la discusión de resultados}. Hoy el capítulo tiene
  resultados y no tiene discusión.
  \item Qué la separa de las secciones anteriores: allí va \emph{qué salió}; aquí, \emph{qué
    significa}. Y sobre todo, qué objetivo queda respondido con cada resultado.
  \item La forma más directa de escribirla, y la que el tribunal puede cotejar de un
  vistazo, es recorrer los seis objetivos específicos y decir de cada uno qué resultado lo
  responde, con qué evidencia y con qué límite. Los que todavía no tengan resultado ---OE4,
  OE6--- se dicen también: un objetivo sin respuesta declarada es mejor que un objetivo que
  el lector no encuentra.
  \item \textbf{Y aquí es donde toca contar lo incómodo}, que es lo que distingue una
  discusión de un resumen: que la recuperación no separa las tres estrategias de
  fragmentación, que la elección se apoya en un efecto sobre la generación que este trabajo
  no ha medido, y que ni la rejilla ni la comparación de modelos exploran el espacio conjunto
  de las dos decisiones.
\end{borrador}

\section{Amenazas a la validez}
\label{secc:amenazas_validez}



\textcolor{borrador2}{De la corrección quedan vivas tres limitaciones que conviene
  declarar. La primera es que las asignaturas de una titulación doble llevan códigos de una
  serie propia, de modo que no se pueden emparejar por código con las de los grados que la
  componen y el emparejamiento tiene que hacerse por el nombre; eso funciona hoy en ciento
  setenta de ciento setenta y ocho casos, pero depende de que la fuente escriba igual la
  misma asignatura en dos planes distintos, y un cambio de redacción lo degradaría en
  silencio. La segunda es que cuando un nombre coincide con varias guías distintas no se
  resuelve la ambigüedad, y esas asignaturas reciben un fragmento que afirma que no tienen
  guía publicada cuando en realidad sí la tienen. La tercera es que una de las cinco
  titulaciones dobles no publica absolutamente nada, de modo que el sistema no puede decir
  qué se estudia en ella.}


\begin{borrador}
  \item \textbf{Dos referencias quedaron colgando al borrar párrafos el 05/08/2026, y hay
    que resolverlas.} El párrafo de arriba empieza por «\emph{De la corrección} quedan vivas
  tres limitaciones», pero ya no hay ninguna corrección mencionada antes: se borró el párrafo
  que contaba que cinco de las doce titulaciones estuvieron en la colección sin una sola
  asignatura hasta el 05/08/2026. O vuelve ese párrafo, o hay que reescribir el arranque de
  este para que no remita a nada.
  \item \textbf{Y más abajo} queda un «CORREGIDO: \emph{el método tachado} se abandonó en
  IT-97» cuyo texto tachado ya no existe.
  \item \textbf{Se borró también el aviso de que las cifras son de otra colección}, y ese sí
  hay que reponerlo en el cuerpo del capítulo y no solo aquí: la comparación de modelos de
  incrustaciones se ejecutó sobre 797 fragmentos y la colección vigente tiene \textbf{1\,334}
  desde IT-16, de modo que las cifras absolutas de la
  Tabla~\ref{tab:comparativa_embeddings} no corresponden al corpus de hoy. Sin esa
  advertencia, el capítulo las presenta como vigentes.
  \item \textbf{Y la distancia ha crecido}: esa tabla va ya dos configuraciones por detrás
  del corpus, porque entre medias entraron los planes de los dobles grados y bajó el máximo
  de fragmento. Lo defendible no es convertir las cifras por regla de tres, sino declararlo
  o volver a ejecutar la comparación. A favor de no rehacerla: el orden de los modelos se
  mantuvo idéntico sobre dos corpus distintos, y eso es lo que sostiene el ADR-0003. En
  contra: ese argumento vale para el orden, no para las cifras absolutas, y son las cifras
  absolutas las que aparecen impresas en la tabla.
  \item \textbf{Fase:} 1 (IT-101). Andamiaje de la amenaza de los dobles grados; la
  evidencia completa, con las seis anomalías de la fuente y sus mediciones, está en
  el DQA-0005 del repositorio\cite{repo_tfg}.
  \item \textbf{Cifras del 06/08/2026:} 528 asignaturas, 288 guías, 8 bloques de salidas,
  \textbf{1\,334} fragmentos (1\,193 de guía, 86 informativos, 33 de plan de estudios y 22
  de salidas). \textbf{599} fragmentos citan una titulación doble, y \textbf{560} de ellos
  son fragmentos de guía reaprovechados: no se duplicó ni un carácter de temario. El
  dataset no cambia al fragmentar más fino ---las 528 asignaturas y las 288 guías son las
  mismas---, solo cambia en cuántos trozos se reparte.
  \item \textbf{Hay que actualizar las cifras del resto del capítulo y del Capítulo~4},
  que siguen siendo las de 797 fragmentos.
  \item \textbf{Y hay un resultado nuevo que este capítulo todavía no cuenta:} sobre los
  1\,334 fragmentos, el acierto por unidad llega a \textbf{1,000 con K${=}$10}, cuando con
  884 se quedaba en 0,995. La pregunta que faltaba era P-008, sobre las salidas
  profesionales de los grados de la rama industrial, que es de agregación y no tenía ningún
  fragmento que la contestara entera. Es una pregunta de cincuenta, así que no conviene
  presentarlo como un salto: lo relevante es que por primera vez no queda ninguna fuera.
\end{borrador}

\begin{borrador}      
  \item \textbf{Fase:} 2 (IT-58). Aquí van todas las amenazas juntas, no solo la de abajo:
  la codificación declarada en falso por la fuente, la cobertura de guías del 82 \% con
  sesgo hacia las titulaciones en implantación, el ECTS ausente, los parámetros de
  fragmentación sin validar y los umbrales de la hipótesis fijados después de IT-28.
\end{borrador}

\begin{borrador}      
  \item \textbf{Fase:} 1 (IT-95) --- amenaza nueva, de \textbf{validez de
    constructo}, y de las más serias que tiene el trabajo. Todo el contenido de
  la colección se obtiene hoy de guías en PDF (288 de 288), y la extracción se
  apoya en una lista de 17 rótulos de sección escrita a mano contra la
  plantilla que la Universidad usaba al observarla. \textit{(Verificado el
    30/07/2026: el código conoce 17 rótulos y espera encontrar 13 en toda guía;
    de esos, solo 2 pasan el filtro de contenido.)}
  \item Si esa plantilla cambia, una sección deja de terminar donde debe: o se queda
  corta, y se pierde contenido, o absorbe la siguiente, y puede arrastrar el
  bloque de profesorado con sus datos personales. \textbf{Ninguno de los dos
    casos se manifiesta como un error}: el sistema sigue respondiendo, solo que
  sobre menos contenido del que cree tener.
  \item \textbf{Actualización:} la amenaza sigue en pie, pero ya está
  \emph{instrumentada}, que no es lo mismo que descartada. El rastreo conserva
  ahora una copia de cada PDF y un verificador comprueba, contra el original,
  que todos los rótulos que la plantilla compone son de los conocidos. Mientras
  esa comprobación salga vacía, las fronteras de sección caen donde el código
  cree; en cuanto devuelva algo, hay que mirarlo antes de fiarse del corpus.
  \item La forma de enunciarla, entonces, es esta: no se puede garantizar que la
  plantilla no cambie —eso no depende del trabajo—, pero sí que un cambio deje
  de pasar inadvertido. Es una diferencia que conviene marcar, porque es la
  misma lección de los cuatro defectos anteriores: lo que no se comprueba no
  
  \item \textbf{Lo que se usa ahora} es la comprobación invertida: en vez de buscar
  rótulos desconocidos, se comprueba qué rótulos \emph{faltan} de los trece
  que la plantilla siempre compone. Un rótulo de más no significa nada; un
  rótulo de menos significa que la plantilla ha cambiado, que es justo lo que
  hay que detectar.
  \item \textbf{Esto es mejor material de defensa que lo que sustituye}, y conviene
  contarlo así: una heurística validada sobre 4 casos aguantó, y se rompió al
  pasarla a 293. Es el mismo patrón que los defectos de la Fase 0 --- las
  pruebas cubren lo que los datos contienen --- pero esta vez ocurrió sobre un
  \emph{verificador}, es decir, sobre la herramienta que existe precisamente
  para que no pase.
  \item Dato útil para la defensa, y para adelantarse a la pregunta obvia sobre
  cuánto se tira: de los \sout{53.969} \textbf{50.715} caracteres de secciones
  que suman las tres guías de prueba se conservan 8.534,
  \sout{alrededor del 15 \%} el \textbf{16,8~\%}. El resto se descarta
  a propósito y con nombre ---evaluación, bibliografía, cláusulas, profesorado,
  metodologías--- y el verificador lo enumera sección por sección, de modo que el
  filtrado se puede revisar en lugar de tener que creérselo.
  \item (El total cambió porque IT-97 movió las fronteras de sección al cambiar cómo
  se detectan los rótulos; los 8.534 conservados son los mismos. Medido de
  nuevo el 30/07/2026 sobre las tres guías de prueba.)
\end{borrador}

\begin{borrador}      
  \item \textbf{Fase:} 1 (IT-95) --- anomalía de la fuente, para el registro de
  calidad de datos y para matizar la cifra de cobertura. Cinco asignaturas de
  planes en implantación publican su guía docente \textbf{con el esqueleto
    completo de la plantilla y las secciones de contenido vacías}.
  \item Se comprobó descargando los seis PDF implicados: los seis se leen
  correctamente, con entre 9.300 y 13.000 caracteres cada uno y los trece
  rótulos de la plantilla en su sitio. En uno de ellos, las líneas
  <<RESUMEN>>, <<DESCRIPCIÓN DE CONTENIDOS>> y <<METODOLOGÍAS DOCENTES>>
  aparecen seguidas, sin nada en medio.
  \item Lo que importa aquí no es el número, sino la distinción: no es un fallo de
  extracción del sistema, es un hueco del origen. El aviso del rastreador decía
  <<PDF ilegible>> y la colección decía <<no ha podido obtenerse>>; las dos
  cosas insinuaban un fallo propio que no existe. Ya está corregido: es
  cuestión de honestidad de la propia colección, que no puede afirmar algo
  falso sobre por qué le falta un dato. Queda registrado en el DQA-0004.
  \item Merece contarse como pareja del defecto de IT-94, porque son el mismo error
  en direcciones opuestas: allí la colección iba a negar una publicación que sí
  existe; aquí iba a atribuirse un fallo que no había cometido. Las dos veces
  el texto era cómodo de escribir y falso, y las dos veces lo que lo destapó
  fue ir a mirar la fuente en lugar de fiarse de lo que el código creía.
  \item Enlaza con la amenaza ya declarada del sesgo de cobertura: las titulaciones
  de implantación reciente no solo concentran las asignaturas sin guía, sino
  también las guías publicadas y vacías, que hasta ahora se contaban como
  cobertura. \textbf{La cifra de cobertura que se publique debería separar
    guías con contenido de guías vacías}, o declarar expresamente que no lo hace.
\end{borrador}

\textcolor{borrador2}{Entre las amenazas a la validez de construcción conviene destacar una por lo
  que enseña sobre el propio método de trabajo. Durante la Fase 0 se dieron tres defectos que la
  batería de pruebas no llegó a detectar y que solo aparecieron al ejecutar el sistema contra el
  conjunto de datos completo. Al revisar el fragmentador apareció un cuarto caso del mismo tipo, con
  una diferencia que merece señalarse: se localizó antes de que los datos llegaran a activarlo. Las
  asignaturas que la fuente publica sin código se identificaban de forma que unas sobrescribían a
  otras, y el encabezado de un fragmento podía terminar atribuyendo el temario a una materia
  distinta. No llegó a ocurrir, porque ninguna de esas asignaturas tiene todavía guía publicada, pero
  tampoco había nada que lo hubiera avisado: el verificador comparaba las claves de los fragmentos y
  no el texto de sus encabezados, de manera que informaba de que el corpus era correcto. La
  conclusión que se extrae no es que las pruebas sobren, sino que su cobertura está limitada por los
  casos que el conjunto de datos contiene en un momento dado, y que un verificador solo protege de
  aquello que efectivamente comprueba.}

\begin{borrador}
  \item \textbf{Comprueba el <<Fase 0>> de este párrafo (07/08/2026).} Los tres defectos que
  las pruebas no detectaron sí son de la Fase~0, pero el cuarto es del fragmentador, y el
  troceado se cuenta ahora en la Fase~1. Tal como está redactado no dice que el cuarto sea de
  la Fase~0, así que puede que no haya que tocar nada; míralo y decide.
  \item Lo mismo con la frase de más arriba que habla de <<el mismo patrón que los defectos
  de la Fase~0>>.
\end{borrador}

\textcolor{borrador2}{La regeneración del corpus al curso siguiente confirmó ese diagnóstico de la
  peor manera posible, y conviene contarlo porque matiza la conclusión anterior. El fragmentador, que
  llevaba meses funcionando sin incidencias, dejó de terminar: no falló ni avisó de nada, sencillamente
  se quedó dando vueltas. El procedimiento que une los fragmentos demasiado cortos daba por hecho que
  repartir de nuevo un par siempre produce una división distinta, cuando el texto solo puede partirse
  por sus fronteras naturales y a veces las únicas disponibles son las que ya lo separaban. Bastó con
  que una guía tuviera dos fragmentos cuya suma excediera el máximo por un solo carácter para que el
  proceso se repitiera indefinidamente. El defecto estaba en el código desde la fase inicial y ningún
  conjunto de datos lo había activado: lo activó el cambio de la fuente, porque las guías del curso
  nuevo llegan en un formato distinto y con ello cambian las fronteras del texto. La lección no es
  solo que las pruebas cubren lo que los datos contienen, sino que la fuente puede modificar el
  comportamiento del sistema sin que se toque una línea de código.}

\textcolor{borrador2}{Ese mismo procedimiento volvió a no terminar por segunda vez, y el
  segundo caso enseña algo distinto del primero. La corrección anterior se apoyaba en un
  razonamiento explícito sobre por qué el proceso acaba: el número de fragmentos solo puede
  disminuir, y como no puede bajar de uno, el número de reintentos está acotado. El
  razonamiento era falso. Repartir de nuevo un par que no cabe junto puede devolver tres
  fragmentos donde había dos, de modo que la cuenta también sube, y alternando uniones que la
  bajan con repartos que la suben el proceso oscila indefinidamente sin llegar a ninguna
  parte. Con los tamaños de fragmento en uso el caso era inalcanzable, y quedó al descubierto
  al hacerlos configurables para poder ensayar valores menores. La corrección consiste en
  descartar todo reparto que produzca más fragmentos de los que había, con lo que la cuenta
  pasa a ser efectivamente decreciente y el razonamiento de terminación se sostiene. Lo que
  merece señalarse no es el defecto en sí, sino que la primera corrección venía acompañada de
  un argumento escrito de por qué no podía repetirse, y ese argumento no cubría todos los
  casos: una demostración informal da una confianza que no siempre está justificada, y la
  prueba de regresión que se añadió entonces no podía detectarlo porque comprobaba
  exactamente el escenario que el argumento sí contemplaba.}
